\chapter{Objetivos del Proyecto}

\section{Objetivo General}

Diseñar, implementar y validar un sistema multi-agente de inteligencia artificial
para la auditoría automatizada de contratos de concesión de infraestructura en el
Perú, que reduzca el tiempo de \textit{screening} inicial de \textbf{320 horas a
menos de 1 hora}, manteniendo cobertura del 100\,\% de cláusulas, precisión
$\geq 85\,\text{\%}$ en la detección de inconsistencias y trazabilidad completa de
cada hallazgo, sin reemplazar el juicio experto del analista en la etapa de
interpretación y decisión.

Este objetivo responde directamente a la pregunta central planteada en la
Sección~2: el sistema actúa como herramienta de \textit{screening} sistemático
que eleva la cobertura, consistencia y transparencia del proceso, liberando al
especialista para concentrarse en el análisis de alto valor.

\section{Objetivos Específicos}

\begin{enumerate}
    \item \textbf{Pipeline de procesamiento documental.}
    Construir un pipeline de ingesta, segmentación e indexación multinivel que
    procese contratos PDF de más de 300 páginas con estructura jerárquica variable,
    incluyendo mecanismo anti-TOC y detección robusta de secciones, cláusulas y
    referencias cruzadas.

    \item \textbf{Arquitectura multi-agente especializada.}
    Implementar tres agentes LLM especializados que operen en paralelo sobre cada
    sección del contrato: el \textit{Jurista} (inconsistencias jurídicas y
    normativas), el \textit{Auditor} (referencias cruzadas y estructura) y el
    \textit{Cronista} (plazos, fechas y obligaciones temporales).

    \item \textbf{RAG híbrido con normativa sectorial.}
    Integrar un módulo de \textit{Retrieval-Augmented Generation} híbrido
    (FAISS\,+\,BM25\,+\,Cohere Reranker) que recupere normativa sectorial peruana
    relevante para contextualizar el análisis de cada sección.

    \item \textbf{Grafo de conocimiento contractual (GraphRAG).}
    Desarrollar un módulo de construcción dinámica de grafos que modele relaciones
    semánticas entre cláusulas, plazos, roles y entidades del contrato, habilitando
    consultas relacionales imposibles con búsqueda vectorial plana.

    \item \textbf{Despliegue en producción.}
    Desplegar el sistema como servicio accesible mediante API REST (Google Cloud
    Run), interfaz web (Next.js) y bot conversacional (Telegram), accesible bajo
    el dominio \texttt{contractia.pe}.

    \item \textbf{Validación experimental y comparativa.}
    Validar el MVP sobre el Contrato de Concesión A (324 páginas, 448 cláusulas)
    y realizar una evaluación comparativa entre Gemini~2.5~Pro y
    Gemini~3.1~Pro~Preview en términos de cobertura, velocidad y calidad de
    hallazgos.
\end{enumerate}

\section{KPIs de Éxito}

Los objetivos son verificables mediante tres indicadores clave de desempeño,
definidos siguiendo la metodología \textbf{SMART} (Specific, Measurable,
Achievable, Relevant, Time-bound) en coordinación con especialistas de
ProInversión:

\begin{landscape}
\begin{table}[H]
\centering
\caption{Indicadores Clave de Desempeño (KPIs) para Validación del Sistema de Revisión de Contratos}
\label{tab:kpis_validacion}
\small % Usar un tamaño de fuente pequeño para que quepa bien

% Usamos 'adjustbox' si la tabla es demasiado ancha, aunque 'small' ya ayuda
\begin{adjustbox}{max width=24cm}
\begin{tabular}{>{\raggedright\arraybackslash}p{5cm} | >{\raggedright\arraybackslash}p{5cm} | >{\raggedright\arraybackslash}p{3.5cm} | >{\raggedright\arraybackslash}p{10cm} | >{\raggedright\arraybackslash}p{6cm} | >{\raggedright\arraybackslash}p{7cm}}
\toprule
\textbf{KPI} & \textbf{Línea Base (Proceso Manual)} & \textbf{Objetivo (Sistema)} & \textbf{Método de Medición} & \textbf{Criterio de Aceptación} & \textbf{Responsable de Validación} \\
\midrule
\multicolumn{6}{l}{\textbf{Eficiencia Operacional}} \\
\midrule
TTR - Time to Review (screening inicial) & 320 horas (rango: 280-400h) & $\le 1$ hora & Timestamp desde carga de documento hasta reporte completo & Promedio sobre $\ge 3$ contratos $< 1$ hora & Equipo técnico + Cronómetro del sistema \\
\midrule
Cobertura de procesamiento & 100\% (manual) & 100\% (automatizado) & Verificación que todas las secciones fueron procesadas & Todas las secciones listadas aparecen en resultados & Especialista Legal \\
\midrule
\multicolumn{6}{l}{\textbf{Precisión de Detección}} \\
\midrule
Precisión (referencias rotas) & N/A & $\ge 90\%$ & Validación experta sobre muestra $n\ge 50$ & $\text{VP}/(\text{VP}+\text{FP}) \ge 0.90$ & Panel de especialistas senior \\
\midrule
Recall (referencias rotas) & N/A & $\ge 85\%$ & Validación experta sobre muestra $n\ge 50$ & $\text{VP}/(\text{VP}+\text{FN}) \ge 0.85$ & Panel de especialistas senior \\
\midrule
Precisión (contradicciones) & N/A & $\ge 85\%$ & Validación experta en muestra $n\ge 30$ & $\text{VP}/(\text{VP}+\text{FP}) \ge 0.85$ & Panel de especialistas senior \\
\midrule
Recall (contradicciones) & N/A & $\ge 80\%$ & Validación experta sobre muestra $n\ge 30$ & $\text{VP}/(\text{VP}+\text{FN}) \ge 0.80$ & Panel de especialistas senior \\
\midrule
Tasa de falsos positivos (global) & N/A & $\le 5\%$ & Revisión manual de los hallazgos en contratos de validación & $\text{FP}/(\text{VP}+\text{FP}) \le 0.05$ & Especialistas por área \\
\midrule
\multicolumn{6}{l}{\textbf{Cobertura y Trazabilidad}} \\
\midrule
Cobertura de indexación de cláusulas & 100\% (manual) & $\ge 97\%$ & Conteo manual en muestra representativa vs conteo del sistema & Ratio $\ge 0.97$ con margen $\pm 3\%$ & Especialista técnico + revisor independiente \\
\midrule
Trazabilidad completa de hallazgos & Variable (documentación ad-hoc) & 100\% & Verificación de campos obligatorios en estructura de datos & 100\% de hallazgos con los 5 campos obligatorios poblados & Script de validación automático \\
\midrule
Auditabilidad de ubicaciones & N/A & 100\% & Verificación manual de muestra aleatoria $n=20$ & Ubicación coincide con contenido citado en 100\% de casos & Auditor independiente \\
\midrule

\multicolumn{6}{p{\dimexpr\linewidth-2\tabcolsep\relax}}{\textbf{Notas sobre la tabla:}} \\ % Línea especial para el encabezado de las notas
\multicolumn{6}{p{\dimexpr\linewidth-2\tabcolsep\relax}}{%
    \begin{itemize}
        \item VP = Verdaderos Positivos, FP = Falsos Positivos, FN = Falsos Negativos.
        \item N/A en línea base indica que el proceso manual actual no captura esta métrica.
        \item Los umbrales de contradicciones son ligeramente menores que referencias rotas reconociendo mayor complejidad de detección semántica.
        \item La cobertura de indexación acepta 97\% (no 100\%) para acomodar casos extremos de formato altamente irregular.
    \end{itemize}
} \\
\bottomrule
\end{tabular}
\end{adjustbox}
\end{table}
\end{landscape}



\newpage

\subsection{KPI-1 — Eficiencia operacional (Tiempo de screening)}

\textbf{Declaración:} Reducir el TTR (\textit{Time-to-Review}), medido desde la
carga del PDF hasta la generación del reporte consolidado, de 320 horas a
\textbf{menos de 1 hora} para contratos de 200--400 páginas, demostrando una
reducción $\geq 99.7\,\text{\%}$ en tiempo de procesamiento inicial.

\textbf{Relevancia:} El 100\,\% de los especialistas entrevistados identificó
el tiempo de revisión como el principal cuello de botella del proceso. Procesar
un contrato en menos de una hora elimina la dependencia de semanas de calendario
para completar el \textit{screening} inicial.

\subsection{KPI-2 — Precisión y cobertura de detección}

\textbf{Declaración:} Lograr precisión $\geq 85\,\text{\%}$ y recall $\geq 85\,\text{\%}$
en detección de inconsistencias críticas (referencias cruzadas rotas,
contradicciones procedimentales, omisiones normativas), con cobertura del
$100\,\text{\%}$ de las cláusulas del contrato, validado mediante comparación con
revisión experta independiente.

\textbf{Relevancia:} Alta precisión previene la sobresaturación del analista
con alertas incorrectas (\textit{alert fatigue}). La cobertura del 100\,\%
garantiza que ninguna cláusula quede sin analizar, a diferencia del proceso
manual que cubre aproximadamente el 60\,\% por muestreo.

\subsection{KPI-3 — Trazabilidad completa}

\textbf{Declaración:} El 100\,\% de los hallazgos reportados debe incluir:
ubicación exacta (sección, cláusula, párrafo), extracto textual de contexto,
tipo de inconsistencia y fundamento normativo cuando corresponda. Ningún
hallazgo puede reportarse sin estos metadatos.

\textbf{Relevancia:} La trazabilidad es el requisito de aceptación institucional
más crítico: sin ella, el sistema no puede ser auditado ni sus resultados
verificados por la Contraloría u otros organismos de control.

\begin{table}[H]
\centering
\caption{Resumen de KPIs de éxito del MVP}
\label{tab:kpis_resumen}
\begin{tabular}{p{4.5cm}p{3.5cm}p{3cm}p{3.5cm}}
\toprule
\textbf{KPI} & \textbf{Línea base (manual)} & \textbf{Objetivo MVP} & \textbf{Resultado obtenido} \\
\midrule
Tiempo de screening (TTR) & 320 horas & $< 1$ hora & $\sim$1h 1min $\checkmark$ \\
Cobertura de cláusulas & $\sim$60\,\% (muestreo) & 100\,\% & 100\,\% $\checkmark$ \\
Precisión semántica & Variable & $\geq 85\,\text{\%}$ & 89\,\% $\checkmark$ \\
Recall (sensibilidad) & Variable & $\geq 85\,\text{\%}$ & 94\,\% $\checkmark$ \\
Trazabilidad de hallazgos & Baja / manual & 100\,\% & 100\,\% $\checkmark$ \\
\bottomrule
\end{tabular}
\end{table}